% Part: sets-functions-relations
% Chapter: sets
% Section: russells-paradox

\documentclass[../../../include/open-logic-section]{subfiles}


\begin{document}

\olfileid{sfr}{set}{rus}
\olsection{رسل دا تضاد}

عنصراں نال تعیّن دا اصول علامتی لکھت $\Setabs{x}{\phi(x)}$ دا جواز دیندا
اے، جہڑی اوہناں $x$ دا \emph{اوہی} سیٹ دسدی اے جنہاں لئی~$\phi(x)$
ہووے۔ پر عنصراں نال تعیّن دا اصول \emph{اصل وچ} صرف ایس خیال دا جواز
دیندا اے: \emph{جے} اک ایہو جہا سیٹ موجود ہووے جہدے رکن سارے اوہ تے صرف
اوہ ہون جیہڑے $\phi$ نیں، \emph{تاں} ایہو جہا سیٹ اکو اے۔ ہور لفظاں وچ:
کِسے~$\phi$ نوں طے کرن مگروں، سیٹ $\Setabs{x}{\phi(x)}$ یکتا اے،
\emph{جے اوہ موجود ہووے}۔

پر ایہ شرط بڑی اہم اے!{} بنیادی گل ایہ اے کہ ہر خاصیت توں
\emph{سیٹ نہیں بنایا جا سکدا}۔ یعنی کجھ خاصیتاں سیٹ دی تعریف
\emph{نہیں} کردیاں۔ جے ساریاں کردیاں، تاں اسی صاف تضاداں وچ پھس جاواں
گے۔ ایس دی سب توں مشہور مثال رسل دا تضاد اے۔

سیٹ ہور سیٹاں دے !!{element}s ہو سکدے نیں---مثال لئی، سیٹ~$A$ دا ذیلی
سیٹاں دا سیٹ، سیٹاں توں ای بنیا ہوندا اے۔ سو ایہ پچھنا یا کھوجنا بجا اے
کہ کوئی سیٹ دوجے سیٹ دا !!a{element} اے یا نہیں۔ کی کوئی سیٹ اپنے آپ دا
رکن ہو سکدا اے؟ سیٹ دے خیال وچ کوئی گل ایس نوں روکندی نہیں لگدی۔ مثال
لئی، جے \emph{سارے} سیٹ شےواں دا اک اکٹھ بناندے نیں، تاں کوئی ایہ سوچ
سکدا اے کہ اوہناں نوں اکو سیٹ وچ اکٹھا کیتا جا سکدا اے---سارے سیٹاں دا
سیٹ۔ تے اوہ، آپ اک سیٹ ہون پاروں، سارے سیٹاں دے سیٹ دا !!a{element}
ہووے گا۔

رسل دا تضاد اُدوں پیدا ہوندا اے جدوں اسی اپنے آپ نوں !!a{element} نہ
رکھن دی خاصیت، یعنی \emph{اپنے آپ دا رکن نہ ہون} دی خاصیت، اُتے غور کردے
آں۔ جے اسی منّ لئیے کہ اوہناں سارے سیٹاں دا سیٹ موجود اے جو اپنے آپ نوں
!!a{element} نہیں رکھدے، تاں کی ہووے گا؟ کی
\[
R = \Setabs{x}{x \notin x}
\]
موجود اے؟ گل ایہ اے کہ اسی ثابت کر سکدے آں کہ ایہ موجود نہیں۔

\begin{thm}[رسل دا تضاد]\ollabel{thm:russells-paradox}
	کوئی سیٹ $R = \Setabs{x}{x \notin x}$ موجود نہیں۔
\end{thm}

\begin{proof}
جے $R = \Setabs{x}{x \notin x}$ موجود ہووے، تاں
$R \in R$ اُدوں تے صرف اُدوں جدوں $R \notin R$، تے ایہ اک تضاد اے۔
\end{proof}

\begin{tagblock}{novice}
\begin{explain}
آؤ ایس ثبوت نوں ہولی ہولی ویکھیئے۔ جے $R$ موجود اے، تاں ایہ پچھنا بجا اے
کہ $R \in R$ اے یا نہیں۔ منّ لو کہ واقعی $R \in R$ اے۔ ہُن $R$ دی
تعریف اوہناں سارے سیٹاں دے سیٹ دے طور تے کیتی گئی سی جو اپنے آپ دے
!!{element}s نہیں۔ سو، جے $R \in R$، تاں $R$ آپ اوہ خاصیت نہیں رکھدا
جس نال $R$ دی تعریف کیتی گئی اے۔ پر صرف اوہ سیٹ $R$ وچ نیں جو ایہ
خاصیت رکھدے نیں، ایس لئی $R$، $R$ دا !!a{element} نہیں ہو سکدا، یعنی
$R \notin R$۔ پر $R$، $R$ دا !!a{element} ہووے وی تے نہ وی ہووے، ایہ
نہیں ہو سکدا؛ سو ساڈے کول اک تضاد اے۔

کیوں جو ایہ فرض کرنا کہ $R \in R$ تضاد ول لے جاندا اے، ایس لئی
$R \notin R$ اے۔ پر ایہ وی تضاد ول لے جاندا اے!{} کیوں جو جے $R \notin R$
ہووے، تاں $R$ آپ اوہ خاصیت رکھدا اے جس نال $R$ دی تعریف کیتی گئی اے،
تے ایس لئی $R$، $R$ دا !!a{element} ہووے گا، بالکل اوہناں ہور سیٹاں
وانگ جو اپنے آپ دے رکن نہیں۔ تے فیر، ایہ نہیں ہو سکدا کہ اوہ $R$ دا
!!a{element} نہ وی ہووے تے ہووے وی۔
\end{explain}
\end{tagblock}

\begin{digress}
اسی سیٹاں دا ایہو جہا نظریہ کِویں بنائیے جو رسل دے تضاد وچ نہ پھسے، یعنی
ایہ \emph{بے مطابقت} دعویٰ نہ کرے کہ $R = \Setabs{x}{x \notin x}$
موجود اے؟ ایس لئی سانوں مسلّمے مقرر کرنے پَین گے، جو ایہ دسن لئی بالکل
درست شرطاں دین کہ سیٹ کدوں موجود ہوندے نیں (تے کدوں نہیں ہوندے)۔

ایس باب وچ سیٹاں دے نظریے دا جیہڑا خاکہ دتا گیا اے، اوہ ایہ کم نہیں کردا۔
ایہ \emph{واقعی سادہ فہم} نظریہ اے۔ ایہ صرف دسدا اے کہ سیٹ عنصراں نال
تعیّن دے اصول اُتے چلدے نیں تے، جے تہاڈے کول کجھ سیٹ ہون، تاں تسیں اوہناں
دا اتحاد، اشتراک وغیرہ بنا سکدے او۔ سیٹاں دے نظریے نوں ایس توں زیادہ
سخت منطقی بنیاد اُتے اساریا جا سکدا اے۔ \oliflabeldef{cumul:::part}{اوہ
سختی حصہ \olref[cumul][][]{part} لئی چھڈی جائے گی۔ ہُن لئی، اسی سادہ فہم
ڈھنگ نال اگے ودھاں گے تے احتیاط نال تضاداں توں بچن دی کوشش کراں گے۔}{}
\end{digress}


\end{document}
